% =============================================================================
% Abstrak Grafis dan Graphical Abstract -- satu halaman
% Gambar akhir disarankan memiliki resolusi sekurang-kurangnya 300 dpi.
% =============================================================================

\ipbprelimheading{ABSTRAK GRAFIS}

\begin{center}
\ifdefempty{\BerkasAbstrakGrafis}{%
  \fbox{\parbox[c][5.7cm][c]{0.8\textwidth}{\centering
    Sisipkan sebuah gambar yang secara visual merangkum temuan utama}}%
}{%
  \includegraphics[width=0.8\textwidth,height=5.7cm,keepaspectratio]{\BerkasAbstrakGrafis}%
}

\vspace{0.3cm}

{\bfseries\parbox{9cm}{\centering\JudulAbstrakGrafis}\par}
\end{center}

\smallskip

\begingroup
\setlength{\parindent}{1.27cm}%
Abstrak grafis/\textit{graphical abstract} terdiri atas sebuah gambar yang
secara visual merangkum temuan utama dari tugas akhir. Abstrak grafis ini
disusun secara singkat dan berdiri sendiri, menyampaikan tujuan dan hasil
penelitian, serta disajikan dengan cara yang menarik secara visual dengan
resolusi gambar minimal 300 dpi. Abstrak grafis/\textit{graphical abstract}
ditulis dalam satu halaman. Judul diletakkan di bawah abstrak
grafis/\textit{graphical abstract}. Mahasiswa wajib mendeklarasikan penggunaan
AI jika digunakan dalam pembuatan abstrak grafis dan perhatikan penggunaan
elemen berhak cipta (lihat Suplemen 6).
\par\endgroup

\begin{otherlanguage}{english}
\itshape
\ipbprelimsubheading{GRAPHICAL ABSTRACT}

\begin{center}

\ifdefempty{\BerkasAbstrakGrafisInggris}{%
  \fbox{\parbox[c][5.7cm][c]{0.8\textwidth}{\centering
    Insert a single image that visually summarizes the main findings}}%
}{%
  \includegraphics[width=0.8\textwidth,height=5.7cm,keepaspectratio]{\BerkasAbstrakGrafisInggris}%
}

\vspace{0.3cm}

{\bfseries\parbox{9cm}{\centering\JudulAbstrakGrafisInggris}\par}
\end{center}
\end{otherlanguage}
